\section{Numerical parity and statistical validity}\label{sec:parity}

The empirical claim of this paper is intentionally narrower than the
package surface. \statspai{} registers \RegistryTotal{} public
callables, and the census behind the claim is reported once, here: \RegistryCertified{}
are \code{certified} and \RegistryValidated{} \code{validated}
(together \RegistryCertifiedValidated{}), \RegistryApiStable{} are
\code{api\_stable}, and \RegistryExperimental{} are experimental. The
word ``validated'' in this submission means that a function is in the
\code{certified} or \code{validated} registry tier, and the tiers are
earned in exactly two ways: \code{certified} by a cross-language
comparison with a named reference on identical bytes, \code{validated}
by known-truth recovery or replication of published numbers. Monte
Carlo coverage rows and documented disagreements (T4) are recorded
against a function but never earn a tier on their own, so a disclosure
cannot be read as a demonstration. All \RegistryStableHandwritten{}
hand-written stable symbols carry parity, reference, or API-contract
evidence; the remaining \RegistryAutoUnbacked{} stable auto-registered
symbols are API-stable but not parity-backed, a denominator revisited
in Section~\ref{sec:discussion}. Every \code{certified} symbol is
linked in the registry to a named external reference implementation and
the tolerance it was held to --- an \proglang{R} or \proglang{Stata}
parity module for all but a handful whose canonical implementation is a
\proglang{Python} package (\pkg{econml}, \pkg{DoubleML}, \pkg{pyblp},
\pkg{obp}); those exceptions are enumerated in the package
(\code{PYTHON\_REFERENCE\_ROWS}) and asserted by a test.
Every \code{validated} symbol is linked to at least one known-truth or
published-replication artifact; the validation-evidence audit in the archive fails on any
certified/validated symbol whose registry entry lacks such an
attachment, so the mapping from the \RParityModuleCount{} modules to
the \RegistryCertifiedValidated{} certified/validated symbols is
enumerable rather than asserted.

Two properties of that census are easy to misread, so both are stated
here rather than left to be recomputed. First, the certified and
validated tiers answer different questions and are not summed in this
paper's coverage language. \ParityCrossLanguage{} symbols were compared
against a named \proglang{R} or \proglang{Stata} implementation on
identical inputs; \ParityInternalEvidence{} recover a known population
parameter on a deterministic DGP or reproduce published numbers on a
calibrated replica. Only the first is a parity claim. The second is the
strongest evidence available for a method with no reference
implementation, and reporting a combined figure would let it borrow the
first one's authority. Second, the all-registered denominator dilutes
coverage: of \RegistryTotal{} registered callables, \ParityClassTotal{}
are result or exception classes that can never carry a parity grade and
\ParityInfraTotal{} are infrastructure that renders tables, draws plots,
builds agent schemas or loads data. Against the \ParityEstimatorTotal{}
estimator callables that remain, cross-language coverage is
\ParityEstimatorCrossLanguage{} functions
(\ParityEstimatorCrossLanguagePct{}). \code{sp.parity\_summary()} returns
this split at call time, so a reader never has to reconstruct the
denominator from a printed total.

The registry tier and the parity grade are derived from one artifact,
the committed parity index, through a single grade-to-tier function; a
test asserts that the two agree for every stable symbol.

The validation design is a twelve-estimator suite with four evidence
tracks: numerical parity (Track A), Monte Carlo coverage against known
DGPs (Track B, separated into materialized \(B_{1000}\) rows and
estimator-specific slow-suite \(B_s\) rows), computational scaling (Track C,
Section~\ref{sec:performance}), and agent-facing mechanical checks
(Track D, Section~\ref{sec:agentbench}).
This is representative, not exhaustive: it covers the causal-inference
core behind the paper's statistical evidence claims, not every
auxiliary helper in the registry.

Because the paper uses several graded vocabularies, Table~\ref{tab:terminology}
fixes them before any evidence is read. The practical
reading rule is as follows. A \code{certified} or
\code{validated} label is local to the named function, data contract,
and evidence kind recorded in the registry; it does not promote nearby
helpers or optional backends, and only a T2 row is a parity claim.
Conversely,
\code{api\_stable} means that the call signature and output shape are
stable enough to support user code, not that the estimator has the
paper-level statistical evidence required for a validation claim.

\begin{table}[tbp]
\centering\small
\begin{tabular}{p{0.24\linewidth}p{0.70\linewidth}}
\toprule
Vocabulary & Meaning \\
\midrule
Registry tiers & \code{certified} / \code{validated} / \code{api\_stable} /
\code{experimental}: the per-function evidence label queryable at call
time (Section~\ref{sec:registry}). \\
Evidence grades T1--T5 & Kind of numerical evidence behind a row:
known-truth recovery (T1), same-byte cross-language parity (T2),
seed-replicated equivalence of a stochastic estimator within a stated
margin (T3), documented
convention or identification boundary (T4), registered but not
validated here (T5) (Section~\ref{sec:validation-tiers}). \\
Tracks A--D & The four evidence tracks of the validation suite: parity,
coverage, performance, agent-facing mechanical checks. \\
\(B_{1000}\), \(B_s\) & Monte Carlo draw counts: materialized
1{,}000-draw rows versus the slower estimator-specific suite. \\
Reviewer paths & The three reproduction routes of
Section~\ref{sec:reproducibility}: the licence-free \proglang{Python}
path, the pinned-\proglang{R} path, and the \proglang{Stata} path. \\
\bottomrule
\end{tabular}
\caption{Terminology used by the validation sections. Each vocabulary is
orthogonal: a \code{certified} function can carry T2 and T4 rows on
different fixtures, and any row can be rebuilt on one or more reviewer
paths.}
\label{tab:terminology}
\end{table}

\subsection{Twelve-estimator validation suite}\label{sec:decathlon}

\begin{table}[!ht]
\centering\small
\begingroup
\setlength{\tabcolsep}{3pt}
\begin{tabular}{@{}rL{0.17\linewidth}L{0.26\linewidth}L{0.24\linewidth}L{0.15\linewidth}l@{}}
\toprule
\# & Method & Entry point & Reference & Data / DGP & Tracks \\
\midrule
1  & OLS + HC SE             & \code{sp.regress}         & \code{lm}+\pkg{sandwich}; \code{regress} & Card 1995 & A,\(B\),D \\
2  & 2SLS                    & \code{sp.iv}              & \pkg{AER}; \code{ivregress} & Card 1995 & A,\(B\),D \\
3  & HDFE                    & \code{sp.fast.feols}      & \pkg{fixest}; \code{reghdfe} & HDFE panel & A,\(B\),C,D \\
4  & Callaway--Sant'Anna     & \code{sp.callaway\_santanna} & \pkg{did}; \code{csdid} & \texttt{mpdta} & A,\(B\),C,D \\
5  & Sun--Abraham            & \code{sp.sun\_abraham}    & \pkg{fixest}; \code{eventstudyinteract} & \texttt{mpdta} & A,\(B\),D \\
6  & Bias-corrected RD       & \code{sp.rdrobust}        & \pkg{rdrobust}; \code{rdrobust} & RDsenate & A,\(B\),D \\
7  & RD density test         & \code{sp.rddensity}       & \pkg{rddensity}; \code{rddensity} & RDsenate & A,D \\
8  & Classical SCM           & \code{sp.synth}           & \pkg{Synth}; \code{synth} & Basque; unique DGP & A,C,D \\
9  & Synthetic DID           & \code{sp.sdid}            & \pkg{synthdid}; \code{sdid} & California-99 & A,\(B\),D \\
10 & DML PLR                 & \code{sp.dml}             & \pkg{DoubleML} & Card/401(k) & A,\(B\),C,D \\
11 & Causal forest           & \code{sp.causal\_forest}  & \pkg{grf} & Clean-overlap DGP & A,\(B\),D \\
12 & PSM                     & \code{sp.psm}             & \pkg{MatchIt}; \code{teffects} & NSW-DW & A,D \\
\bottomrule
\end{tabular}
\endgroup
\caption{Twelve-estimator validation suite. In the Reference column the
\proglang{R} package precedes the \proglang{Stata} command. Tracks are A
(parity), \(B\) (materialized 1,000-draw coverage), C (performance), and
D (agent-facing mechanical checks); row 8 dispatches as
\code{sp.synth(method="classic")} and row 10 as
\code{sp.dml(model="plr")}. Data marked with a named extract are
calibrated replicas in Track~A and original extracts in
Table~\ref{tab:headline-parity}. The HDFE member's coverage row
runs through the \code{sp.panel} fixed-effects entry of the same
absorption family. Supplemental Track-B rows for
2x2 DiD and entropy balancing are reported in
Table~\ref{tab:track-b-b1000} but are not members of this suite.}
\label{tab:decathlon}
\end{table}

The reference implementations of Table~\ref{tab:decathlon} are cited
once here: \code{lm} with \pkg{sandwich} \citep{zeileis2020sandwich},
\pkg{AER}'s \code{ivreg} \citep{kleiber2008applied}, \pkg{fixest}
\citep{berge2018efficient}, \pkg{did} \citep{callaway2021difference},
\pkg{rdrobust} \citep{calonico2014robust,calonico2015rdrobust},
\pkg{rddensity} \citep{cattaneo2018manipulation,cattaneo2020simple},
\pkg{Synth} \citep{abadie2011synth}, \pkg{synthdid}
\citep{arkhangelsky2021synthetic}, \pkg{DoubleML}
\citep{bach2022doubleml,bach2024doubleml}, \pkg{grf}
\citep{athey2019generalized}, and \pkg{MatchIt} \citep{ho2011matchit},
with the \proglang{Stata} commands \code{csdid} \citep{rios2022csdid},
\code{sdid} \citep{clarke2024synthetic}, \code{eventstudyinteract}
\citep{sun2021estimating}, and \code{reghdfe}'s Poisson companion
\code{ppmlhdfe} \citep{correia2020fast}.

\subsection{Estimands, identification, and inference contracts}
\label{sec:statistical-contracts}

Numerical agreement does not establish identification. This subsection
therefore states the contract under which each validation-suite row has a
causal interpretation. These are
contracts on the user's design and data, not properties that software can
certify from a successful function call.  In particular, parity with a
reference implementation can establish that two programs compute the same
quantity while both remain inappropriate for a design that violates
exogeneity, parallel trends, continuity, or overlap.

For the linear rows, write \(Y=X\beta+u\).  The OLS coefficient is always a
linear-projection parameter; interpreting it causally additionally requires
the relevant conditional-mean restriction.  With instruments \(Z\), 2SLS
solves the sample analogue of \(E[Z(Y-X\beta)]=0\); a LATE interpretation for
a binary treatment further requires relevance, independence, exclusion, and
monotonicity.  Absorbing fixed effects changes the projection space, not these
identification requirements.  The parity harness specifies the covariance
estimator explicitly---HC1 or the named cluster---rather than relying on the
entry point's non-robust default.

For staggered adoption, the target primitive is

\[
  \operatorname{ATT}(g,t)
  = E\!\left[Y_t(1)-Y_t(0)\mid G=g\right],
\]

identified under no anticipation, a stated never-treated or not-yet-treated
comparison group, conditional parallel trends, overlap, and no interference.
Callaway--Sant'Anna aggregates these group--time effects; Sun--Abraham first
forms cohort-by-relative-time effects and then interaction weights.  A joint
pre-trend test is reported, but failure to reject is not proof of parallel
trends and its power must be considered.

The RD row targets the sharp/fuzzy cutoff effect (or kink contrast) under
continuity, local support, and no precise manipulation; its parity row
uses the native robust bias-corrected interval at the native CCT
MSE-optimal bandwidth (Section~\ref{sec:track-a}).

The remaining rows target nonparametric discontinuities, counterfactual
panels, or effects learned after nuisance estimation.  For RD density the null
is equality of the one-sided density limits at the cutoff; rejecting it is a
design diagnostic, not itself an effect estimate.  Classical SCM and SDID
construct weighted untreated counterfactuals, so pre-treatment fit and the
stability of the untreated outcome process are central.  Their placebo
distributions are design-based diagnostics and should not be read as generic
i.i.d. sampling distributions. SDID uses placebo inference by default
(with bootstrap and jackknife alternatives), while classical SCM offers
in-space placebos.

DML PLR uses an orthogonal, cross-fitted score for the partially linear model
\(Y=\theta D+g(X)+\varepsilon\), \(D=m(X)+v\).  A causal reading of
\(\theta\) requires the corresponding conditional exogeneity and overlap
conditions in addition to nuisance-rate conditions.  The forest row targets
\(\tau(x)=E[Y(1)-Y(0)\mid X=x]\); its ATE/ATT summaries use an AIPW score and
therefore require conditional exchangeability and overlap.  PSM targets ATT
by default and requires conditional exchangeability, common support, and a
credible propensity specification.  Balance and overlap diagnostics remain
part of the analysis even when the reference parity row passes. The validation
path uses five-fold DML cross-fitting, AIPW influence-function intervals for
forest aggregates, and Abadie--Imbens heteroskedastic inference for default
nearest-neighbour PSM.

\subsection{Track A: Numerical parity}\label{sec:track-a}

Track A uses three layers. First, estimators are tested against known
DGP truths whenever a truth exists. Second, \statspai{} and the
canonical \proglang{R} implementation read the same \code{\%.16g} CSV
bytes and write full-precision JSON. Third, when a directly comparable
\proglang{Stata} command exists, the Stata side reads the same CSV and
emits a sibling JSON artifact. The current R-joined harness has
\RParityModuleCount{} modules. Of those, \StataParityModuleCount{} carry a
canonical or audited Stata reference; each of the
\NoStataModuleCount{} without one records a measured, command-level
reason in \code{compare.py::STATA\_SKIP\_REASON}.
Rows whose defaults differ for a documented reason are labelled as
convention gaps rather than as ordinary passes.
Table~\ref{tab:internal-parity} summarizes the source-tree test layers
behind these claims.

\begin{table}[!ht]
\centering\small
\begin{tabular}{p{0.25\linewidth}p{0.22\linewidth}rp{0.36\linewidth}}
\toprule
Layer & Harness & Tests & Role \\
\midrule
Analytical/reference recovery & reference parity & 4{,}285 & Deterministic-DGP recovery, frozen \proglang{R} / \proglang{Stata} reference fixtures on identical bytes, and cross-estimator agreement; includes the 203 tests that pin Stata's \code{vce()} grammar (27 estimators, 79 estimator-by-covariance cells at \(10^{-6}\)) and its \code{test}, \code{lincom} and \code{margins} post-estimation commands \\
Published/replica parity & external parity & 88 & Pinned values for Card, Lee, Basque, NSW, \texttt{mpdta}, paper-level formulas, and the four \pkg{DoubleML} model-class pins \\
Monte Carlo coverage & coverage MC & 32 & Slow-marked headline DGP checks, robustness rows, and one smoke row \\
\bottomrule
\end{tabular}
\caption{Validation test layers available in the source tree used for
this submission. Counts are the tests \code{pytest} collects in each
harness, reported by \code{sp.validation\_report(collect\_tests=True)}
and pinned as floors by \code{tests/test\_jss\_validation\_api.py}; the
Track A modules of Appendix~\ref{app:parity} are a separate,
module-level artifact.}
\label{tab:internal-parity}
\end{table}

\begin{table}[!ht]
\centering\scriptsize
% AUTO-GENERATED by tests/r_parity/compare.py
% Compact main-manuscript Track A snapshot; do not hand edit.
\setlength{\tabcolsep}{2pt}
\begin{tabular}{@{}>{\raggedright\arraybackslash}p{0.195\linewidth}>{\raggedright\arraybackslash}p{0.106\linewidth}>{\raggedright\arraybackslash}p{0.084\linewidth}r@{\hspace{3pt}}r@{\hspace{3pt}}r@{\hspace{4pt}}r@{\hspace{4pt}}>{\raggedright\arraybackslash}p{0.115\linewidth}@{}}
\toprule
Estimator & Statistic & Data & \statspai{} & \proglang{R} ref. & \proglang{Stata} ref. & Max rel. err. & Verdict \\
\midrule
\code{sp.fast.feols} & \(\hat\beta_{x1}\) & 2-way HDFE, \(N{=}10^4\) & 1.999503980 & 1.999503980 & 1.999503980 & \(2.9\times10^{-15}\) & T2; fixest ssc \\
\code{sp.regress} & \(\hat\beta_{\mathrm{educ}}\) & Card replica & 0.109986249 & 0.109986249 & 0.109986249 & \(6.3\times10^{-13}\) & T2 \\
\code{sp.iv} & \(\hat\beta_{\mathrm{educ}}\) & Card replica & 0.141834409 & 0.141834409 & 0.141834409 & \(1.1\times10^{-11}\) & T2 \\
\code{sp.callaway\_santanna} & simple ATT & \texttt{mpdta} replica & -0.032976514 & -0.032976514 & -0.032976514 & \(1.3\times10^{-15}\) & T2 \\
\code{sp.rdrobust} & robust RD, default CCT \(h\) & RDsenate & 7.506502484 & 7.506502484 & 7.506502483 & \(9.4\times10^{-11}\) & T2; native, default \(h\) \\
\code{sp.dml} & \(\hat\theta_{\mathrm{DML}}\) & Card replica & 0.110228260 & 0.110228260 & 0.110228260 & \(3.7\times10^{-15}\) & T2; ddml bridge \\
\code{sp.causal\_forest} & AIPW ATE & clean-overlap DGP & 1.012009376 & 1.011741552 & --- & 0.000265 & T3; seed-replicated, operator exact \\
\code{sp.psm} & ATT & NSW--DW replica & 2277.690721 & 2277.690721 & 2277.690721 & \(1.2\times10^{-15}\) & T2; SE convention \\
\code{sp.synth} (unique) & avg post gap & identified SCM DGP & 2.000000155 & 1.999999999 & 2.000000000 & \(7.8\times10^{-8}\) & T1+T2; identified DGP \\
\code{sp.synth} (classic) & avg post gap & Basque replica & -0.704244148 & -0.687789210 & -0.703950536 & 0.0239 & T4; reference disagreement \\
\code{sp.frontier} & intercept & frontier DGP & 2.071558153 & 2.071558166 & 2.071558168 & \(7.2\times10^{-9}\) & T2 \\
\code{sp.decompose} & Oaxaca gap & Oaxaca DGP & 0.221231026 & 0.221231026 & 0.221231026 & \(6.3\times10^{-16}\) & T2 \\
\bottomrule
\end{tabular}

\caption{Representative Track A cross-language parity rows compiled
from the same JSON artifacts as the supplementary parity ledger. Every
row is inside its registered module-level tolerance in
\code{tests/r\_parity/compare.py}, which is \(10^{-6}\) except where
the verdict says otherwise: the DML row's tighter \(10^{-10}\) budget
reflects its fully deterministic linear-learner pipeline, the forest
row is a single draw from each engine whose interpretation rests on the
seed-replicated comparison of Section~\ref{sec:validation-tiers}, and
the classical-SCM row is the \(5\times10^{-2}\) T4 disclosure. Data
marked ``replica'' are the calibrated redistribution-safe fixtures; the
Card-replica rows use the specification \code{lwage \~{} educ + exper +
expersq + black + south + smsa} with HC1 errors, which is why their
coefficients differ from the shorter original-extract specification of
Table~\ref{tab:headline-parity}.}
\label{tab:track-a-cross-language-snapshot}
\end{table}

Table~\ref{tab:track-a-cross-language-snapshot} is only a compact slice;
the full row-level ledger is generated in the replication archive. Its
fixtures are calibrated, redistribution-safe replicas (the Data column
names them as such); original public extracts are handled separately in
Table~\ref{tab:headline-parity}. Of
the \RParityModuleCount{} R-joined modules, \RParityPassCount{} receive a pass-type verdict, but two
non-T2 Track A rows remain explicitly visible in the parity artifacts:
classical SCM on the calibrated Basque replica is a T4
identification/reference-disagreement disclosure,
and causal forest is a stochastic T3 equivalence claim rather than a
deterministic T2 claim.
A parity row is evidence about a \statspai{} algorithm only if the
\proglang{Python} side ran one, so every module is also classified by
what its \statspai{} side executes, and the classification is not taken
on trust: a call-trace audit
(\code{scripts/trace\_parity\_provenance.py}) runs each module under a
profiler that records every call from \statspai{} into a package other
than the numerical substrate (\pkg{NumPy}, \pkg{SciPy}, \pkg{pandas})
and every subprocess it launches, and a contract test fails if the
registered classification and the trace disagree or if a module has
changed since it was traced. Of the \RParityModuleCount{} modules,
\ParityNativeModuleCount{} run a \statspai{}-native implementation
(\pkg{scikit-learn} appears only as a user-chosen nuisance learner);
\ParityThirdPartyModuleCount{} are served by a third-party
\proglang{Python} estimation library that \statspai{} wraps --- panel
fixed and random effects through \pkg{linearmodels}, ARIMA through
\pkg{statsmodels}, and fixed-effects GLMs through \pkg{pyfixest} ---
and are marked as such in Appendix~\ref{app:parity}, because they
validate the wrapper's argument mapping and conventions rather than an
independent algorithm; \ParityPortModuleCount{} rely on an official
port; and \ParityBackendModuleCount{} call a reference implementation.
That last count was two before release 1.31: until then the
two Honest-DiD modules called \code{backend="honestdid"}, which hands the
computation to the \proglang{R} package itself, so their rows compared
\proglang{R} with \proglang{R}; the RD headline delegated to the
official \pkg{rdrobust} \proglang{Python} port. All three now exercise
native code, and each carries evidence of its own. The native
smoothness FLCI is compared with \pkg{HonestDiD} after the package's
simulated folded-normal quantile is replaced by the exact one (the
shipped package sits \(3\times10^{-4}\) away, which is that simulation
error): the bounds agree to \(1.4\times10^{-8}\) for \(M>0\), and at
\(M=0\), where the interval has a closed form (the GLS linear
extrapolation of the pre-period coefficients), the native solver is
within \(2\times10^{-9}\) of it while \pkg{HonestDiD}'s cone solver is
\(5.7\times10^{-6}\) away. The native relative-magnitudes set is a port
of the Andrews--Roth--Pakes conditional test \citep{andrews2023inference}
inverted over \pkg{HonestDiD}'s \(\theta\) grid; on sixteen cases
(diagonal and correlated covariances, a non-basis target, a single
post-period) it returns the \emph{same} accepted grid points as
\pkg{HonestDiD}, and the least-favourable hybrid, whose first-stage
critical value is simulated, agrees to within one grid step. The native
RD default selector returns \pkg{rdrobust}'s MSE-optimal bandwidth and
robust estimate on the RDsenate bytes at relative error below
\(3\times10^{-13}\), with the port recorded only as a convergence check.
Module 52
separately certifies the native classical-SCM solver on a uniquely
identified DGP. A generated ledger in the replication archive
classifies the remaining T4 row and fails the audit if a methodological
gap appears without an explicit reviewer-risk label and promotion path;
the causal-forest row is governed by the seed-replicated comparison of
Section~\ref{sec:validation-tiers}.

Native \code{sp.sdid} matches \pkg{synthdid}'s ATT at relative error
\(2.6\times10^{-15}\) and native \code{sp.rddensity} matches
\pkg{rddensity}'s p-value at \(3.3\times10^{-11}\) \citep{arkhangelsky2021synthetic,cattaneo2020simple}.

\paragraph{Headline original-data rows.}
The Track~A fixtures above are calibrated, redistribution-safe
replicas; the complementary original-data ledger runs \statspai{} and
the canonical \proglang{R} reference on the \emph{original} public
extracts shipped by \proglang{R} packages. It currently holds
\OrigParityModuleCount{} modules and \OrigParityRowCount{} statistic
rows, including six epidemiology-side modules that reproduce the
NHEFS worked examples of Hern\'an and Robins' \emph{What If} (IPW,
g-formula, g-estimation, outcome regression, survival, and E-values).
Table~\ref{tab:headline-parity} shows the most reviewer-relevant
cases; the full report is generated by
\code{tests/orig\_parity/compare\_orig.py}, and the printed table is
itself generated from those artifacts. The \code{did::mpdta}
simple-ATT row makes the replica/original boundary concrete: the
calibrated replica used by Track~A module 04 returns \(-0.0330\) on all
three sides, while the original bytes return \(-0.03995\); both are
same-byte parity facts, and neither number is quoted against the other
data set anywhere in this paper.

\begin{table}[!ht]
\centering\scriptsize
\resizebox{\textwidth}{!}{%
% AUTO-GENERATED by replication/scripts/generate_headline_parity_table.py
% Cells come from tests/orig_parity/results/*.json; do not hand edit.
\begin{tabular}{@{}L{0.24\linewidth}L{0.19\linewidth}rrrL{0.23\linewidth}@{}}
\toprule
Extract & Statistic & Published & \statspai{} & \proglang{R} & Note \\
\midrule
\code{wooldridge::card} & OLS educ & $0.075$ & $0.07401$ & $0.07401$ & rel(sp,R) \(7.9\times10^{-15}\) \\
\code{wooldridge::card} & 2SLS educ & $0.132$ & $0.1323$ & $0.1323$ & rel(sp,R) \(1.6\times10^{-11}\) \\
\code{did::mpdta} & CS simple ATT & -- & $-0.03995$ & $-0.03995$ & rel(sp,R) \(2.3\times10^{-14}\); no printed vignette anchor \\
\code{did::mpdta} & CS dynamic overall ATT & $-0.0772$ & $-0.07724$ & $-0.07724$ & rel(sp,R) \(1.6\times10^{-14}\); vignette prints $-0.0772$ \\
\code{Synth::basque} & ADH avg gap & $-0.855$ & $-0.8946$ & $-0.8944$ & rel(sp,R) \(1.8\times10^{-4}\); Stata \code{synth} $-0.8943$ on the same bytes \\
\code{rdrobust::RDsenate} & robust RD jump & -- & $7.507$ & $7.507$ & rel(sp,R) \(3.5\times10^{-16}\); \code{bwselect="cct"} path \\
\code{causalsens::lalonde.psid} & adjusted OLS ATT & $700$ & $751.9$ & $751.9$ & rel(sp,R) \(2.2\times10^{-13}\); public composite extract \\
NHEFS (\emph{What If}) & IP-weighted ATE & $3.4$ & $3.441$ & $3.441$ & rel(sp,R) \(1.3\times10^{-11}\) \\
NHEFS (\emph{What If}) & g-formula ATE & $3.5$ & $3.463$ & $3.463$ & rel(sp,R) \(7.4\times10^{-12}\) \\
\bottomrule
\end{tabular}
}
\caption{Selected original-data parity rows on public extracts,
generated from \code{tests/orig\_parity/results/}. All \proglang{R}
references read the same CSV bytes as \statspai{}. The \emph{Published}
column is a heterogeneous anchor --- an original-paper table value
(Card, Basque), a printed-vignette value (the \code{did} dynamic ATT),
a textbook program output (NHEFS), or a rounded composite-extract value
(LaLonde) --- shown for context, not as a parity target; the parity
claim is the \statspai{}-versus-\proglang{R} pair in the \emph{Note}
column. The \code{did::mpdta} simple-ATT row has no printed vignette
anchor, so its Published cell is empty by design.}
\label{tab:headline-parity}
\end{table}

\subsection{Evidence tiers and hard cases}\label{sec:validation-tiers}

\paragraph{Evidence is attached to configurations.}
A tier is attached to a function, but every artifact exercises one
configuration: a method variant, an inference option, a data condition,
a code path, and an output. The Track~B DML row is the plain example:
the suite member is the partially linear model, and until this revision
the only materialised coverage row ran the interactive model, which says
nothing about the other's intervals (a PLR row is now included,
Table~\ref{tab:track-b-b1000}). For the twelve suite estimators the
package therefore records, per artifact, the configuration it exercises,
and \code{sp.validation\_scope(result)} reads the configuration of a
fitted object and returns the artifacts that match it exactly
(``covered'' requires a T1 or T2 row; ``stochastic only'' means only
T3, T4, or coverage rows match; otherwise ``not covered''), together
with the artifacts that differ in one dimension. The map is deliberately
narrower than the registry --- an option it does not list is reported as
not covered --- and it found a gap in this paper's own first example:
\code{sp.iv}'s default classical 2SLS covariance, printed in
Listing~\ref{lst:card-output}, had no deterministic reference row
(Track~A pins the HC1 variant). It is now pinned against
\pkg{AER}\code{::ivreg} on the original Card extract, with classical,
HC1, and clustered errors in just- and over-identified models, at
\(8\times10^{-11}\).

\paragraph{Evidence grades.}
The grades are those of Table~\ref{tab:terminology}; T5 marks a function
registered but not validated in this paper. Two rows need more than a
tolerance to read correctly.

\paragraph{Causal forest: which Monte Carlo error?}
For a stochastic estimator two uncertainties must be kept apart:
sampling variation across datasets, which the reported standard error of
the ATE estimates, and the algorithm's own Monte Carlo variation across
seeds on \emph{fixed} data. ``Agreement within combined Monte Carlo
error'' is a statement about the second, so it cannot be tested by
dividing the gap between one \statspai{} fit and one \pkg{grf} fit by
their sampling standard errors --- which is how an earlier version of
this row was graded. The forest evidence now separates three tasks.
\emph{Truth recovery and inference} are Track~B's (the AIPW coverage row
below). \emph{Estimand agreement} holds by construction: both engines
report the doubly-robust AIPW average \citep{athey2019generalized}, and
given the same forest the AIPW operator agrees with \pkg{grf}'s to
machine precision. \emph{Algorithmic agreement} is tested by refitting
each engine under 50 seeds (20 at 8{,}000 trees) on fixed data, at
500, 2{,}000, and 8{,}000 trees, on two datasets: the reference-parity
fixture and the Track~A module-13 design, the latter run after the
equivalence margin was fixed. Seeds are spaced \(10^5\) apart on both
sides, because \pkg{grf} forests grown from consecutive seeds share most
of their random draws: a first run with consecutive \proglang{R} seeds
understated \pkg{grf}'s seed-to-seed spread several-fold and made
\statspai{}'s forest look several times noisier than it is.
Table~\ref{tab:forest-seed-mc} reports the result. At every tree count
on both datasets a two-one-sided test establishes equivalence within
\(\delta = 0.1\) sampling standard errors --- a difference that moves a
95\% interval endpoint by at most 5\% of its half-width --- with an
upper bound of at most 0.056. The two engines are equally noisy: their
seed-to-seed standard deviations differ by a factor between 0.78 and
1.48 and fall with the number of trees, from 13--18\% of the sampling
standard error at 500 trees to 6--8\% at the default 2{,}000. The claim
is equivalence, not identity: the engines are independent
implementations, and at 8{,}000 trees, where Monte Carlo error is
smallest, their seed means differ by up to 0.025 sampling standard
errors (\(|z| \le 2.6\)). The single-fit Track~A row is one draw from
each distribution, not a parity fact on its own.
Artifacts: \path{tests/reference_parity/test_grf_seed_mc_equivalence.py}
and its fixtures.

\begin{table}[!ht]
\centering\footnotesize
% AUTO-GENERATED by replication/scripts/generate_forest_seed_table.py
\begingroup\setlength{\tabcolsep}{3pt}
\begin{tabular}{@{}lrrrrrrrl@{}}
\toprule
Data & Trees & Seeds & Gap & MC SE & Gap/SE\(_s\) & SD\(_{\mathrm{MC}}\) (\statspai{}) & SD\(_{\mathrm{MC}}\) (\pkg{grf}) & Equiv. \\
\midrule
Reference fixture & 500 & 50/50 & -0.00028 & \(2.0\times10^{-3}\) & -0.004 & \(1.0\times10^{-2}\) & \(1.0\times10^{-2}\) & yes \\
 & 2000 & 50/50 & +0.00135 & \(9.2\times10^{-4}\) & +0.018 & \(4.4\times10^{-3}\) & \(4.9\times10^{-3}\) & yes \\
 & 8000 & 20/20 & +0.00139 & \(6.2\times10^{-4}\) & +0.019 & \(2.3\times10^{-3}\) & \(1.5\times10^{-3}\) & yes \\
Module-13 design & 500 & 50/50 & +0.00003 & \(9.8\times10^{-4}\) & +0.001 & \(4.7\times10^{-3}\) & \(5.1\times10^{-3}\) & yes \\
 & 2000 & 50/50 & +0.00015 & \(4.6\times10^{-4}\) & +0.004 & \(2.2\times10^{-3}\) & \(2.4\times10^{-3}\) & yes \\
 & 8000 & 20/20 & -0.00003 & \(3.7\times10^{-4}\) & -0.001 & \(1.0\times10^{-3}\) & \(1.3\times10^{-3}\) & yes \\
\bottomrule
\end{tabular}
\endgroup

\caption{Seed-replicated comparison of \code{sp.causal\_forest} and
\pkg{grf}~2.6.1 on fixed data (AIPW ATE; ATT in the archive).
``Gap'' is the difference of seed means; ``MC SE'' its combined
algorithmic Monte Carlo standard error; ``Gap/SE\(_s\)'' the gap in units
of \pkg{grf}'s reported sampling standard error; ``MC SD'' the seed-to-seed
standard deviation of each engine. Equivalence (TOST, 5\%) against
\(\delta = 0.1\,\mathrm{SE}_s\) holds in every row. Generated by
\code{replication/scripts/generate\_forest\_seed\_table.py}.}
\label{tab:forest-seed-mc}
\end{table}

\paragraph{Synthetic control: one specification, two data sets.}\label{sec:scm-identification}
The classical-SCM row is a T4 disclosure because the outer \(V\)
optimization is non-convex and donor weights need not be unique --- and
how much that matters turns out to be data-dependent, which the ledger
reports rather than averages away. On the \emph{calibrated Basque
replica} of Track~A module 07 (the same ADH special-predictor
specification and CSV bytes on all three sides), native \statspai{}
tracks Stata \code{synth} for the average post-treatment gap at
relative error \(4.2\times10^{-4}\), while R \pkg{Synth} and Stata
\code{synth} differ from each other by \(2.3\times10^{-2}\). On the
\emph{original} \code{Synth::basque} bytes under the same
specification, the three implementations essentially reconvene:
\statspai{} \(-0.8946\), \proglang{R} \(-0.8944\), \proglang{Stata}
\(-0.8943\), pairwise within \(4\times10^{-4}\)
(Table~\ref{tab:headline-parity}) --- still above the \(10^{-6}\)
strict-parity gate, consistent with solver tolerance on a non-convex
outer problem, but far from the replica's split. The same specification
therefore admits a near-common optimum on one data set and measurably
different local optima on another; a validation ledger that pooled the
two facts into one tolerance would hide exactly the sensitivity a
practitioner needs to know about. The solver is
separately certified on an identified SCM DGP (Track~A module 52): when
the treated unit is exactly a unique convex combination of donors,
\statspai{}, \pkg{Synth}, and Stata \code{synth} agree on the average
post-treatment gap at relative error \(7.8\times10^{-8}\) and on the
donor weights at \(5.6\times10^{-7}\), the tolerance of the constrained
quadratic solver. Thus the
paper claims solver correctness on identified problems and documents
the replica's reference-disagreement/local-optimum problem without
treating it as a strict R/Synth parity pass.

Native \code{sp.augsynth} now ports \pkg{augsynth}'s Ridge+SCM
convention \citep{benmichael2021augmented} (ATT at \(7.9\times10^{-6}\)
on the Basque fixture).
Native \code{sp.gsynth} now ports \pkg{gsynth}'s two-way factor
convention \citep{xu2017generalized}, matching it to machine precision.

\subsection{Track B: Monte Carlo coverage}\label{sec:track-b}

Track B checks 95\% interval coverage on known DGPs, and reports with
each rate what is needed to explain it: the estimator's bias, its Monte
Carlo SD, the mean reported standard error, their ratio, and the
interval length (Table~\ref{tab:track-b-b1000}). The artifact holds
twelve \(B=1{,}000\) rows; ten correspond to suite members (OLS, 2SLS,
fixed effects through \code{sp.panel}, Callaway--Sant'Anna,
Sun--Abraham, sharp RD, SDID, DML PLR, and the causal forest), and 2x2
DiD, entropy balancing, and the interactive DML model are supplemental.
Three suite members have no coverage row by construction: the RD density
test is a diagnostic, classical SCM has no analytic interval, and PSM's
inference evidence is the Abadie--Imbens SE parity row.

Most rows are calibrated in the direct sense: the mean standard error is
within 3\% of the Monte Carlo SD and coverage sits in the 99\% Wilson
band around 0.95. The rows that are not are explained rather than
graded. \emph{Sharp RD} (0.934) and \emph{SDID} (0.928) have calibrated
standard errors (ratios 0.98 and 0.99), so the shortfall is in the shape
of the sampling distribution, not the SE. For RD this is the reference
procedure's own finite-sample behaviour: on the first 200 draws R
\pkg{rdrobust}~4.0.0 returns the same intervals as \statspai{} to
\(6\times10^{-12}\) and covers in exactly the same draws
(\path{tests/coverage_monte_carlo/mechanisms/rd_drawwise_reference.py}).
\emph{DML PLR} under-covers clearly (0.883), and the diagnostics say why:
its bias is half a Monte Carlo SD while the SE is nearly calibrated
(0.88). \code{sp.dml} returns \pkg{DoubleML}'s numbers to \(10^{-16}\) on
the same folds and learners, so this is the procedure with the default
learner, not the code. Varying one thing at a time on the same design
(\path{tests/coverage_monte_carlo/mechanisms/dml_plr_learners.py},
\(B=500\)) locates it in the nuisance fit: with the true nuisance
functions coverage is 0.952 and the bias vanishes; with a random forest
or a spline Lasso it is 0.940 and 0.938; with the default gradient
boosting it is 0.892 at \(n=500\) and 0.936 at \(n=2{,}000\), where the
bias is still 0.41 SD. DML's orthogonal score makes nuisance error
second order, not zero, and the default learner's regularisation bias on
this design shrinks too slowly to be negligible at these sample sizes;
the result is stated in \code{sp.dml}'s documentation and in the
configuration map. \emph{Entropy balancing} over-covered at 1.000 in the
previous revision because its standard error held the weights fixed and
was twice the Monte Carlo SD; with the M-estimation variance it is now
0.945 with a ratio of 0.97 (Section~\ref{sec:parity-gaps}).
The \emph{causal-forest} row is new in its present form. It previously
over-covered (0.977, standard-error ratio 1.13) because its entry point
reported a separate AIPW estimator rather than the forest's
(Section~\ref{sec:attachment-audit}); recomputed on the forest's own
doubly-robust average at the default 2{,}000 trees it covers 0.959 with
a ratio of 1.04, and a 300-draw check at 300 trees gives 0.973
(\path{tests/coverage_monte_carlo/mechanisms/forest_trees.py}).
Table~\ref{tab:track-b-robustness} records stress cases as failure-mode
evidence rather than calibration results.

% AUTO-GENERATED by replication/scripts/generate_track_b_tables.py
\begin{table}[!ht]
\centering\small
\begingroup\setlength{\tabcolsep}{3pt}
\begin{tabular}{@{}L{0.25\linewidth}L{0.21\linewidth}rrrrrr@{}}
\toprule
Estimator & DGP & Cover. & Bias & MC SD & Mean SE & SE/SD & CI len. \\
\midrule
\code{sp.regress} (HC1) & RCT with covariates & 0.952 & -0.0028 & 0.0707 & 0.0708 & 1.00 & 0.277 \\
\code{sp.regress} 2x2 DiD & 2-period homogeneous DiD & 0.955 & +0.0017 & 0.0974 & 0.0999 & 1.03 & 0.392 \\
\code{sp.ivreg} (HC1) & strong binary-Z IV & 0.962 & -0.0041 & 0.3687 & 0.3627 & 0.98 & 1.422 \\
\code{sp.callaway\_santanna} & homogeneous staggered ATT & 0.947 & -0.0007 & 0.1080 & 0.1090 & 1.01 & 0.427 \\
\code{sp.sun\_abraham} & homogeneous staggered ATT & 0.950 & -0.0007 & 0.1214 & 0.1234 & 1.02 & 0.484 \\
\code{sp.panel} two-way FE & known-coefficient FE panel & 0.948 & -0.0024 & 0.1245 & 0.1233 & 0.99 & 0.486 \\
\code{sp.rdrobust} sharp & known-jump RD, robust CI & 0.934 & -0.0061 & 0.1222 & 0.1201 & 0.98 & 0.471 \\
\code{sp.sdid} (placebo SE) & factor-model panel, 1 treated & 0.928 & +0.0003 & 0.3095 & 0.3071 & 0.99 & 1.204 \\
\code{sp.ebalance} & CIA, ATT (M-estimation SE) & 0.945 & -0.0014 & 0.0799 & 0.0774 & 0.97 & 0.303 \\
\code{sp.dml} (IRM) & binary-\(D\) CIA, IRM ATE & 0.968 & -0.0074 & 0.1582 & 0.1653 & 1.04 & 0.648 \\
\code{sp.dml} (PLR) & continuous-\(D\) PLR, default learners & 0.883 & -0.0261 & 0.0533 & 0.0472 & 0.88 & 0.185 \\
\code{sp.causal\_forest} (AIPW) & binary-\(D\) CIA, population ATE & 0.959 & +0.0095 & 0.0934 & 0.0968 & 1.04 & 0.379 \\
\bottomrule
\end{tabular}
\endgroup
\caption{Track B materialized coverage artifacts, all at \(B=1{,}000\). Next to the coverage rate: the estimator's bias, its Monte Carlo SD, the mean reported SE, their ratio (SE calibration: 1 is exact, above 1 conservative), and the mean interval length. The 99\% Wilson reference band around nominal 0.95 is approximately \([0.935,0.967]\). \code{sp.ivreg} is the same estimator as the \code{sp.iv} alias of Table~\ref{tab:decathlon}. The two DML rows are distinct configurations: the PLR row is the suite member; the IRM row (binary treatment, through \code{sp.causal\_question}) is reported separately and does not stand in for it. The forest row reports the fitted forest's own doubly-robust ATE at the default 2{,}000 trees.}
\label{tab:track-b-b1000}
\end{table}

\begin{table}[!ht]
\centering\small
\begin{tabular}{@{}L{0.25\linewidth}L{0.27\linewidth}rL{0.30\linewidth}@{}}
\toprule
Estimator & Stressor & Coverage & Interpretation \\
\midrule
\code{sp.ivreg} & weak instrument, median \(F\approx2.5\) (\(B=1{,}000\)) & 0.882 & Expected under-coverage; preflight warns and suggests LIML/AR. \\
\code{sp.callaway\_santanna} & heterogeneous timing and magnitude (\(B=1{,}000\)) & 0.946 & Post-fix influence-function scaling remains near nominal. \\
\code{sp.causal\_forest} & severe overlap loss (\(B=300\)) & 0.900 & Under-covers as propensities approach 0 and 1; audit flags overlap before interpretation. \\
\bottomrule
\end{tabular}
\caption{Track B robustness rows. These are descriptive failure-mode checks, not pass/fail calibration claims.}
\label{tab:track-b-robustness}
\end{table}


% Table 11 floats to the top of the next page; without this the
% size/power paragraph split around it and a page opened with the
% fragment "reaches at least 0.82 at the largest effect".
\Needspace*{9\baselineskip}
The size/power sweep closes a weakness of coverage-only reporting: an
interval that never rejects can cover well without being useful. At the
tested null, 2x2 DiD is conservative (size 0.024) and sharp RD rejects
slightly too often (0.076 at \(B=500\)), consistent with its coverage;
the other rows are within Monte Carlo error of 0.05, including entropy
balancing, whose size was 0.000 under the old standard error. Each curve
is monotone on its own effect grid and reaches at least 0.90 at the
largest effect; effect grids and DGPs differ, so this is not a
cross-estimator ranking. DML, the forest, and SDID have no packaged
power sweep, and the paper makes no power claim for them. Full rejection
probabilities are in
\code{tests/coverage\_monte\_carlo/results\_b1000/size\_power\_b1000.json}.

\subsection{Cross-estimator parity}\label{sec:cross-estimator}

The test suite also asserts equivalence across estimators that should
collapse to the same estimand on a given DGP. For example, the
Callaway--Sant'Anna, Sun--Abraham, multi-period DiD, and BJS/imputation
implementations must agree on a homogeneous staggered-DiD DGP to within
a multiple of \(\sqrt{\mathrm{SE}_1^2+\mathrm{SE}_2^2}\). Two estimators fitted
to the same data are positively correlated, so this bound is loose: it is
a screen for gross disagreement, not an equivalence test. It is still a
useful adversarial check, because a shared bug can pass one-reference
parity but fail cross-estimator agreement.

\subsection{High-\texorpdfstring{$N$}{N} analytical parity and pinned fixtures}\label{sec:high-n}

High-\(N\) fixtures tighten selected known-truth checks to two standard
errors at \(N=5{,}000\) and \(N=8{,}000\). Pinned fixtures lock the paper
replica values in Table~\ref{tab:headline-parity} to four decimals.
Honest-DiD is pinned directly to the closed-form breakdown value in
\citet{rambachan2023more}, and the DML/HDFE slices check influence
functions and residual-regression identities independently of any
single reference implementation.

\paragraph{Event-study paths.}
Value comparison is the wrong test for an event-study \emph{path},
because implementations differ in what each half of the path is
differenced against \citep{roth2026interpreting}.
\code{sp.event\_study\_convention()} records that choice per estimator,
and running it closed a defect: \code{sp.did\_imputation} documented
Stata's \code{pretrends(k)} but shipped a residual construction whose
leads equal the symmetric benchmark times $N_0/N$
\citep{li2025benchmarking}, so a violated assumption read as satisfied.
The default now reproduces Stata's leads and their standard errors to
$1.2\times10^{-12}$ and $8\times10^{-13}$.

\subsection{Adjudicating a disagreement}\label{sec:adjudication}

A parity ledger is only as credible as its procedure for the rows that
do \emph{not} agree, so the procedure is fixed in advance and every
registered tolerance carries the outcome. A disagreement on identical
bytes is assumed to be a \statspai{} defect until the first diverging
intermediate quantity -- design matrix, weights, residuals, sandwich
meat -- has been located. If the divergence is a documented choice on
both sides (a degrees-of-freedom divisor, a small-sample cluster
correction, an analytic versus influence-function variance), the
default follows the method's own authors, the alternative is exposed as
a keyword, and the gap is registered with its mechanism and bounded by
the tolerance. If the reference itself is the problem or the solution
is not unique, the row is graded T4 and must carry independent
evidence (a second reference, an analytic identity, or a uniquely
identified DGP). If exact agreement is impossible by construction
(forest randomization, bootstrap), the row is graded T3, which requires
refitting on fixed data under many seeds and an equivalence margin
stated against the sampling uncertainty; a single draw per side is
reported but does not earn the grade. A gap that
cannot be located is not called parity: it is graded ``unexplained''
in the tolerance audit and reported here as an open item. Two contract
tests enforce the procedure mechanically: every \proglang{R}-joined
standard-error row of every passing module must sit inside its
module's registered budget, and every Stata standard-error row that
does not must be registered with a reconstruction of the Stata number.
In release \StatsPAIVersion{} all \RParityModuleCount{} modules pass the
first test, and the second holds three entries, each reproduced from
\statspai{}'s own quantities: the group-average SE of Stata
\code{csdid} (the fixed-cohort-share aggregation of the joint influence
functions, $10^{-14}$; module 04), the PLIV degrees-of-freedom factor
($N/(N-K)$) of Stata's \code{ddml} double-machine-learning command
(module 71), and the per-horizon $K$ of Stata's \code{lpdid} local
projections DiD command (module 83). The Sun--Abraham row is the procedure's worked
example: it was carried as an open 0.7--2.2\% item across two releases
rather than absorbed into a wider tolerance, and the bisection that
closed it in 1.24.0 found a defect on the \statspai{} side
(Table~\ref{tab:correctness-history}).

\subsection{Auditing the attachment, not only the number}
\label{sec:attachment-audit}

A parity harness checks that two implementations produce the same
number. It does not check the claim that sits one level above: that the
function a reader calls is the function the harness ran. Track A modules
are written against one entry point, while a package of this size
exposes dispatchers, thin wrappers and convenience names that are
\emph{said} to reach the same estimator core. Every such name inherits a
grade it did not earn on its own, and inheritance is where an
unfalsifiable claim can hide -- the numbers in the ledger stay correct
while the set of functions they vouch for quietly grows.

Auditing those attachments in release 1.26.0 found that, of eleven
registered alias attachments, seven cited a module that never called the
aliased name. Each surviving attachment now names a test that runs the
alias against the canonical entry point on the module's own bytes. One
did not survive: \code{sp.wooldridge\_did}, credited with the ETWFE
module as an alias of \code{sp.etwfe}, returns an ATT \(6.1\%\) away on
that module's bytes because it is a different estimator; the alias was
withdrawn and the function demoted to the grade its own evidence
supports.

The same failure mode appeared in this revision in a coverage row.
Track~B's causal-forest row ran through
\code{sp.causal\_question()}, whose forest design fitted the forest and
then reported a separate cross-fit AIPW estimator with its own
gradient-boosting nuisances: refitting with 30, 300, and 2{,}000 trees
gave bit-identical results. It now reports the forest's own
doubly-robust average, identical to
\code{cf.average\_treatment\_effect()}, and the row was recomputed.

\subsection{Correctness history and residual risk}\label{sec:parity-gaps}

Every output-changing fix is pinned by a regression test and listed,
with its version, in Appendix~\ref{app:corrections}. Two entries explain
the design better than the rest. The 1.23.0 fix shows what a reference
comparison catches that unit tests do not: \code{sp.did(..., weights=)}
never forwarded the weights to a staggered branch and silently returned
the unweighted estimand, which only \code{did::att\_gt(weightsname=)}
on identical bytes exposed. The 1.25.0 fix shows the limit of a parity
ledger: the causal-forest aggregation was wrong for a
\emph{continuous} treatment, a path no Track~A or Track~B row
exercised, and it surfaced in a worked example. A ledger validates the
configurations it runs and is silent about the rest.

The current release controls that residual risk in three ways rather
than by pointing to the number of past fixes. First, the evidence for a
fitted configuration is queryable: given a result,
\code{sp.validation\_scope} reports which artifacts exercise exactly the variant, inference option,
data condition, and code path of that fit, and reports ``not covered''
rather than inferring coverage from a neighbouring configuration
(Section~\ref{sec:validation-tiers}). Second, the release is frozen:
the paper describes tag \code{v\StatsPAIVersion{}}, the archive carries
the source snapshot and every artifact hash, and changes after the tag
cannot alter a number printed here. Third, every recorded artifact can
be rebuilt from its entry point (Table~\ref{tab:reproduction}).

\subsubsection{Documented default-convention divergences}\label{sec:convention-gaps}

The remaining disclosed boundaries are classical SCM's non-unique
\(V,W\) solutions on the calibrated Basque replica
(Section~\ref{sec:scm-identification}) and the \proglang{Stata} bridge,
which is migration evidence with frozen JSON and do-files and requires a
licence only for re-execution.
